% from ega4-2-fr.tex
\oldpage[IV-2]{113}

Passons à la seconde partie de la proposition. L'existence et l'unicité de
$v$ résultent de ce que pour tout $\alpha$, $(i_\alpha)_*(u)$ est l'unique
homomorphisme rendant commutatif le diagramme
\[
\xymatrix{
\mathcal F\ar[r]^u\ar[d] & \mathcal F'\ar[d]\\
(i_\alpha)_*(\mathcal F|U_\alpha)
\ar[r]_{(i_\alpha)_*(u)} &
(i_\alpha)_*(\mathcal F'|U_\alpha)
}
\]
de ce qu'il y a un unique homomorphisme $w$ rendant commutatifs tous les
diagrammes
\[
\xymatrix{
(i_\alpha)_*(\mathcal F|U_\alpha)
\ar[r]^{(i_\alpha)_*(u)}\ar[d] &
(i_\alpha)_*(\mathcal F'|U_\alpha)\ar[d]\\
\mathcal H^0_{X/Z}(\mathcal F)\ar[r]_w &
\mathcal H^0_{X/Z}(\mathcal F')
}
\]
et enfin de ce que $\mathcal F'$ et
$\mathcal H^0_{X/Z}(\mathcal F')$ s'identifient canoniquement par hypothèse.

Reste à voir que si les supports de $\operatorname{Ker}(u)$ et de
$\operatorname{Coker}(u)$ sont contenus dans $Z$, $v$ est un isomorphisme.
Il suffit de voir que pour tout ouvert noethérien $V$, l'homomorphisme
correspondant
\[
\Gamma(V,\mathcal H^0_{X/Z}(\mathcal F))
\longrightarrow\Gamma(V,\mathcal F')
\]
est alors un isomorphisme. Or, si une section
$t\in\Gamma(V,\mathcal H^0_{X/Z}(\mathcal F))$ a pour image $0$ dans
$\Gamma(V,\mathcal F')$, notons que pour un indice $\alpha$, on a
$t\in\Gamma(V\cap U_\alpha,\mathcal F)$, et en vertu de l'hypothèse sur
$u$, on a $t_y=0$ pour tout $y\in V\cap(X-Z)$; il y a par suite un ouvert
contenant $V\cap(X-Z)$ tel que la restriction de $t$ à cet ouvert soit
nulle, donc par définition $t$ est l'élément $0$ de
$\Gamma(V,\mathcal H^0_{X/Z}(\mathcal F))$.

Prouvons maintenant que toute section $s'\in\Gamma(V,\mathcal F')$ est
l'image d'une section de $\mathcal H^0_{X/Z}(\mathcal F)$ au-dessus de $V$.
Par hypothèse, pour tout $x\in V\cap(X-Z)$ il existe une section $s^{(x)}$
de $\mathcal F$ au-dessus d'un voisinage ouvert $W^{(x)}$ de $x$ dans $X$,
dont l'image par $u$ est $s'|W^{(x)}$; $s'|W^{(x)}$ est donc aussi l'image
par $v$ de la section $t^{(x)}$ de $\mathcal H^0_{X/Z}(\mathcal F)$, image
canonique de $s^{(x)}$. En outre, comme on a vu que $v$ est injectif, les
restrictions de $t^{(x)}$ et $t^{(x')}$ à $W^{(x)}\cap W^{(x')}$ sont
identiques pour deux points quelconques $x$, $x'$ de $V\cap(X-Z)$; les
$t^{(x)}$ sont par suite restrictions d'une même section $t$ de
$\mathcal H^0_{X/Z}(\mathcal F)$ au-dessus d'un voisinage ouvert $U$ de
$(X-Z)\cap V$. Mais comme $\mathcal H^0_{X/Z}(\mathcal F)$ est $Z$-clos,
$t$ se prolonge d'une seule manière en une section de
$\mathcal H^0_{X/Z}(\mathcal F)$ au-dessus de $V$, dont l'image par $v$ a
même restriction que $s'$ à $U$, et coïncide donc avec $s'$ pour la même
raison.

On dit que $\mathcal H^0_{X/Z}(\mathcal F)$ est la $Z$-clôture de
$\mathcal F$.

